The study is closed. These seven topics distinguish completed tests from aspects untested at closure, collected under ``What Remains Unproven'' (Section~\ref{sec:conclusiones}):
\begin{enumerate}
\item \textbf{Burn severity and speed of runoff concentration.} Sentinel-2 dNBR was calculated using scenes from 23 December 2025 and 6 February 2026: \SI{82}{\percent} of the catchments shows a burn signal and \SI{10}{\percent} has moderate-high or high severity. That estimate was incorporated into the balance with assumed runoff coefficients by severity and into the model of \citet{gartner2014}, whose median mixture volume became \SI{148000}{m^3} (Section~\ref{sec:res-incendio}). Runoff increases and the time of concentration of each subcatchment by severity remain untested; the hydrograph rise time to peak remains assumed (\SI{10}{min}).
\item \textbf{Higher-resolution terrain on the campus.} The micro-simulations and m36 reproduced entry through the clearing into the courtyard; m36 calculates \SI{0.17}{m}. The actual geometry of embankments, paths and access points remains unverified by a local LiDAR or photogrammetric survey, which was unavailable in this study.
\item \textbf{Separating rainwater from mud.} Untested: a passive tracer marking the fluid from the hydrographs would allow the perimeter to be evaluated using only that contribution, without counting urban rainfall runoff.
\item \textbf{Porous buildings and coarse load.} Mud retention in homes and additional resistance were tested in a simplified way (v25--v27; porosity verified in z26). Their calibration using damage information by dwelling remains untested; stones, logs and vehicles were not simulated as transported load.
\item \textbf{Independent validation.} Untested for fits made after the held-out test. The facilities test was performed and had an unfavorable outcome (19/49); the fits were not compared with Sentinel-1 radar, official deposit surveys or high-water marks.
\item \textbf{Concentration as a random variable.} Concentration was treated as uncertain through $C_v=0.20+0.30X$ with $X\sim\mathrm{Beta}(3,2)$ (mode 0.40 and mean 0.38), and propagated using five Gaussian nodes (q1--q5). An expected area of \SI{169}{ha} and an ECI between 71 and 73 were obtained (Section~\ref{sec:res-posterior}). Alternative distributions and the estimation of inundation probabilities and extreme percentiles remain untested. Three precautions remain. The distribution, constructed from literature ranges and event evidence, is a belief distribution rather than the actual statistical distribution of sediment, so its results are conditional on it and comparison with alternatives (uniform, Beta(2,2)) was not tested. Quadrature estimates moments well but not inundation probabilities or extreme percentiles, because wetting and drying introduce discontinuities, which can be handled by interpolating between simulations the threshold concentration at which each cell exceeds a given thickness, or with more simulations (Latin hypercube sampling). Varying concentration between simulations represents uncertainty in a parameter, which is a different level of modeling from the physical evolution of concentration along the flow path.
\item \textbf{Evolution of concentration.} Erosion, deposition and incorporation of debris change concentration along the flow path. Untested: an input sediment hydrograph (a more concentrated front and a more dilute tail) and a concentration tracer transported with the flow would allow yield stress and viscosity to vary with it. This distinguishes the final deposit from the maximum flow extent, which is the only one reported at closure.
\end{enumerate}
